% =====================================================================
% METHODOLOGICAL ESSAY — Eight routes to one conclusion
% Long-form essay in the cosmiCave register.
% Drafted 2026-05-29 as the reward artefact after the handoff was committed.
% =====================================================================

\documentclass[11pt,a4paper]{article}

\usepackage[margin=1.1in]{geometry}
\usepackage{amsmath, amssymb}
\usepackage{mathrsfs}
\usepackage[dvipsnames]{xcolor}
\usepackage{hyperref}

\hypersetup{
  colorlinks=true,
  linkcolor=blue!50!black,
  urlcolor=blue!50!black
}

\setlength{\parskip}{0.6em}
\setlength{\parindent}{0pt}

\title{\Large Eight routes to one conclusion\\[0.4em]
\large A methodological essay on foundational physics,\\
\large with no-horizon-formation as the worked example}
\author{Daryl Janzen\\
\small Department of Physics and Engineering Physics\\
\small University of Saskatchewan}
\date{}

\begin{document}
\maketitle

\section*{1. The cave}

In Plato's allegory, the prisoners are not stupid. They are not lazy. They are not failing to do the work they have been trained to do. The work they have been trained to do is the careful study of shadows on a wall, and they do that work with great care. They develop sophisticated taxonomies of shadow-shapes, predictive theories of shadow-motion, and rigorous tests by which competing shadow-theories can be compared and the better one accepted. They build, in their cave, a science.

What the prisoners cannot do --- not because they lack intelligence but because of where they are sitting --- is recognise that the entire enterprise is the study of shadows of something else. The cave makes the shadows look like the primary objects. The shadow-science makes them look like the only objects. To turn around and look at what is casting the shadows requires not better shadow-science but a different posture. It requires leaving the chair.

This essay is about a methodological move in foundational physics that is, structurally, the leaving of the chair. I will describe the move, show it in operation through a single worked example (the no-horizon-formation conclusion in classical general relativity), and argue that what makes the move robust is not the technical sophistication of any one derivation but the architectural feature that the same conclusion is established along eight structurally distinct routes that mutually entail one another. The robustness is not a property of any one route. It is a property of the constellation.

I will also argue that this architectural feature is the right methodological response to a specific institutional dynamic --- one in which careful technical work that bears on a foundational question can be deflected through mechanisms that look, to the field, like normal-science engagement. A single careful argument can be deflected by a single careful deflection. Eight careful arguments operating on eight different terrains cannot.

The essay does not assume technical training in general relativity. It does assume a willingness to follow careful reasoning. The mathematical content is held to what is needed for the argument to be visible. Where formal references are needed, they point to the papers of the programme this essay synthesises~\cite{JanzenBHcausality, JanzenCircle, JanzenSlicing, JanzenGroupoid, JanzenCRcosmology, JanzenFramework}; readers who want the formal proofs can find them there.

\section*{2. What conceptual logic from minimal premises is}

The prisoners' shadow-science is not failed by lack of intelligence. It is failed by where the prisoners sit. The methodological move I want to describe is, accordingly, not a refinement of the shadow-science; it is a different posture toward the work, one available to anyone willing to take it.

The posture has two aspects that operate together as one stance, rather than as two stances in tension. The first is that speculation runs as far as it runs, without restraint, without checking whether the destination will be comfortable for the field, without softening for politeness or strategic positioning. The second is that the speculation, however far it runs, is grounded at every step in what the existing formalism actually says --- not in what the field has been reading it as saying, not in canonical received summaries, not in interpretive habits that have become invisible through repetition, but in the careful reading of the formalism's actual content under premises the formalism itself supplies.

These two aspects are not in tension. The first is the absence of constraint mechanisms on where investigation can go. The second is the demand that wherever it goes is reached by steps grounded in primary-source material at the strength that material actually carries. Together they describe an investigation that is completely open in its scope and obsessively conservative in its grounding, with no managed balance between the two: the openness is total because the grounding is total, and conversely.

What this produces, applied to a question in foundational physics, is \emph{conceptual logic from minimal premises}. The premises are minimal in the specific sense that they are what the existing formalism already supplies --- not new postulates introduced to make the argument work, but the structural features of the formalism that any working practitioner already accepts. The logic is conceptual in the sense that the work is in reading the formalism carefully, tracking what its standard definitions actually imply when applied without the interpretive habits the field has accumulated around them. The conclusion is forced rather than proposed, because what comes out is what the formalism, carefully read, contains, and not what the investigator wishes to advance.

A move of this kind has three load-bearing features that distinguish it from the standard mode in foundational physics.

First: no new formalism is introduced. The investigation does not propose a new field, a new constant, a new equation, a new Lagrangian, or an ontology added on top of the existing one. The investigation reads what is already there, at the level of structural content, without addition.

Second: the conclusion is not a thesis to be debated. A thesis is something the investigator advances, that the field then accepts or rejects according to its standards of argument. The conclusion of a conceptual-logic-from-minimal-premises derivation is not in this category. It is what the formalism, under the minimal premises, exhibits as structural content. The investigator who reports the conclusion is reporting what is in the formalism, not advancing a position.

Third: the same conclusion, when it is genuinely structural to the formalism, can typically be exhibited along multiple structurally distinct routes. The derivation that establishes the conclusion via one route can be redone via another route, on different terrain, against different counter-arguments, with different load-bearing premises. The structural-redundancy architecture this produces is not redundancy in the loose sense of repetition; it is multiple independent demonstrations that the same structural content is present, each operating in a different formal register.

These three features --- no new formalism, conclusion forced rather than proposed, exhibited along structurally distinct routes --- are what make the move robust against the specific class of failures by which careful arguments on foundational questions are deflected in practice. I defer the discussion of that robustness to Section 10, after the worked example has made the move visible in operation.

What the move requires, to be carried out, is the absence of the mechanisms that ordinarily constrain investigation in the field. Three such mechanisms are worth naming explicitly, because their removal is constitutive of the posture and not incidental to it. \emph{Ego}, in the sense of the investigator's attachment to producing the kind of result that will reflect well on the investigator. \emph{Prejudice}, in the sense of the interpretive habits the field has trained the investigator into, operating invisibly until something forces them to surface. \emph{Agenda}, in the sense of any commitment the investigator has to producing one outcome rather than another. Each of these, if active during the investigation, produces a specific failure mode in which the formalism gets read through the constraint rather than at its actual content. The investigator operating under ego will fill in expected structural elements of an argument by reaching for what the field will recognise, rather than by following the formalism's actual content. The investigator operating under prejudice will read the formalism's standard definitions through the interpretive habits the field has trained them in, missing the careful reading the formalism rewards. The investigator operating under agenda will identify in advance the conclusion the investigation is to reach, and then construct the route to it. All three produce conclusions that look like results but are interpretive constructions; the discipline is the removal of these mechanisms, so that what comes out of the investigation is what the formalism actually says.

This is the methodological move the essay articulates. It is not a procedure that can be applied as a checklist; it is a posture that the investigator either takes or does not take. The worked example, in the next several sections, exhibits the posture in operation on a specific question in classical general relativity. The exhibition is meant to do the work of articulation that abstract description cannot.

\section*{3. The worked example: black hole formation}

The example I will work through in detail is one such case: the question of whether a black hole's event horizon forms in finite time, in finite-time sense from any observer external to the gravitational collapse that supposedly produces it.

I should say at the outset what the standard answer is, what the field has been carrying, and what the careful re-derivation reveals. The standard answer is that yes, of course, black holes form. We see them in the sky --- not literally see them, but we see their gravitational effects on stars and gas, we see the gravitational waves released when two of them merge, we see the supermassive black hole at the centre of our galaxy through the orbits of stars near it. The question is read, in the field's training, as one whose answer is observationally settled.

But the question I want to ask is not whether the gravitational phenomena we observe are real. They are real. They are exactly as the standard theory of general relativity describes them. The question is what the gravitational phenomena are caused by, and specifically whether the underlying cause includes the finite-time presence of event horizons --- the closed null surfaces that mark a boundary across which causal influence cannot travel outward.

This is a finer question than it sounds. The standard reading treats it as the same question as ``do black holes exist?'' and reads the question as observationally settled. The careful re-derivation will show that the two questions are not the same question, and that the second is open even though the first is closed.

Let me begin with the minimal premises.

\section*{4. The minimal premises}

The premises I will work from are the following. None is controversial. All are accepted by every working relativist.

\textbf{Premise 1: Causal locality.} Physical influence in classical general relativity propagates along the manifold's causal structure. No signal travels faster than light. Equivalently: if event $A$ is to influence event $B$, then $A$ must lie in the causal past of $B$.

\textbf{Premise 2: The global causal definition of the event horizon.} The event horizon $H^+$ of a black-hole spacetime is, by the standard global causal definition, the boundary of the causal past of future null infinity:
\[ H^+ \;=\; \partial J^-(\mathscr{I}^+). \]
In ordinary language: an event lies inside the horizon if no causal signal from it can ever reach an observer infinitely far in the future. The horizon is, definitionally, the surface that separates events that can in principle send signals to infinity from events that cannot.

\textbf{Premise 3: The asymptotic structure of null infinity.} Future null infinity $\mathscr{I}^+$ is the limiting surface that null geodesics (paths of light rays) approach as their affine parameter tends to infinity. It is not a finite-distance surface within the spacetime; it is the structure that gives meaning to the phrase ``in the limit of arbitrarily long times.''

\textbf{Premise 4: Standard structure of null geodesic congruences.} The local geometry of null geodesics in classical general relativity is the standard one: null geodesics have well-defined tangent vectors, satisfy the geodesic equation, and admit the standard formalism of expansion, shear, and rotation.

These premises are not novel. They are the working content of any general-relativity course. Every result I will derive from them is a result that any competent relativist could in principle derive, given enough time and the right question. The work of the derivation is conceptual: it is in reading what the premises, carefully tracked, actually imply.

\section*{5. The first route: the global theorem}

Consider a star collapsing gravitationally. The collapse begins at some finite exterior time, on the worldline of an exterior observer. The standard reading says that at some finite later time, the collapsing matter has shrunk past its Schwarzschild radius, and an event horizon has formed around it. Inside the horizon, the matter continues to collapse to a singularity. From the moment of horizon formation onward, the interior is causally disconnected from the exterior, and an observer outside the horizon can see only the exterior region.

But now ask: where is the horizon-formation event? The horizon-formation event is the moment at which the collapsing matter's surface crosses the surface that will become the horizon. By the standard definition (Premise 2), the horizon is $\partial J^-(\mathscr{I}^+)$ --- the boundary of the causal past of future null infinity. The horizon-formation event therefore lies on this boundary in the strict sense: it is an event $E$ such that no causal curve from $E$ reaches $\mathscr{I}^+$, but every event in the future of $E$ that lies arbitrarily close to $E$ on the outside of the horizon does have causal curves reaching $\mathscr{I}^+$.

By Premise 1 (causal locality), an event $A$ that can influence the horizon-formation event $E$ must lie in the causal past of $E$. But the horizon-formation event $E$ lies on the boundary of $J^-(\mathscr{I}^+)$, and no causal curve from $E$ reaches $\mathscr{I}^+$. So no signal from $E$ can ever reach any external observer. The horizon-formation event is forever causally disconnected from every event in the external universe.

This is the metric singularity theorem of the programme's first paper~\cite{JanzenBHcausality}, stated informally. What it establishes is that the horizon-formation event, properly so-called, is not a finite-time event from the perspective of any external observer. It is an asymptotic feature of the spacetime --- a feature of the infinite-time null boundary at $\mathscr{I}^+$, not a feature that any finite-time observer can locate in their causal past.

The standard reading, which says that ``at some finite later time, the horizon has formed,'' is incompatible with the premises. The conclusion is forced.

\section*{6. The second route: the asymptotic alignment of null directions}

The first route operates globally: it tracks the horizon-formation event through the causal structure of the entire spacetime. A second route operates locally: it tracks the directions in which null geodesics propagate near where a horizon would form.

Consider a null geodesic that just barely escapes the gravitational collapse --- one whose initial direction is at the boundary of being captured. Such a geodesic, in the limit of barely escaping, takes an arbitrarily long time to reach a distant observer. As the collapse proceeds, the marginally-escaping null directions become more and more nearly tangent to the surface that would, asymptotically, be the horizon.

Now apply Premise 4 (the standard structure of null geodesic congruences) carefully. A horizon, in the geometry of general relativity, is generated by null geodesics whose directions align in a specific way. The horizon's null generators must, asymptotically, lie precisely along the direction that would be tangent to the surface in question. But ``asymptotically'' here means as the affine parameter along the generator tends to infinity --- which, by Premise 3, is the limit at $\mathscr{I}^+$.

What this means concretely is: the alignment of horizon generators along the surface that the standard reading calls ``the horizon'' is achieved only in the limit of infinite time. At any finite time, the would-be generators are not yet aligned along the would-be horizon; they are still propagating outward, just very nearly tangent to it. The horizon as a finite-time geometric structure requires the generators to be already aligned along it; the structure of null geodesic congruences forbids this alignment from happening in finite time.

This is the asymptotic alignment lemma of the programme's first paper, stated informally. It is a local result --- it concerns the geometry of null directions near the would-be horizon, not the global causal structure of the spacetime --- but it reaches the same conclusion as the global route. The horizon is not a finite-time structure. It is the limiting structure at the infinite-time null boundary.

Two distinct routes. Two different terrains: global causal structure in the first, local geometry of null congruences in the second. Two distinct sets of load-bearing premises (Premise 2 was load-bearing in route 1; Premises 3 and 4 are load-bearing in route 2). One conclusion.

\section*{7. The third route: what an observer actually finds}

The first two routes are technical. The third route is what one finds if one tries to trace, concretely, what an observer attempting to chart the horizon-formation event would experience. It is the methodological move at its most accessible: not a theorem but a thought experiment.

Imagine an astronaut falling into a black hole, carrying a magic flashlight that emits a flash every second on the astronaut's own clock. The flashes propagate outward, and a distant observer records the times at which they arrive.

At first, the flashes arrive at the distant observer at slightly delayed intervals: as the astronaut falls deeper into the gravitational well, the flashes have farther to climb out of, and each takes a little longer to reach the distant observer. The intervals are seconds-plus-a-bit on the distant observer's clock for each second on the astronaut's clock. Standard general relativity describes this with precision; it is the gravitational redshift.

But as the astronaut approaches what the standard reading calls the horizon, something specific happens to the intervals. The delay grows. Not by a constant factor, not by a finite multiplier, but without bound. The flash emitted at the astronaut's local time $t = T - 1$ second (one second before crossing the horizon, in the astronaut's frame) reaches the distant observer at some time. The flash emitted at the astronaut's local time $t = T - 0.1$ second reaches the distant observer at a much later time. The flash emitted at the astronaut's local time $t = T - 0.001$ second reaches the distant observer at a still much later time. The flash emitted at the astronaut's local time $t = T$ --- the moment of horizon crossing, in the astronaut's frame --- never reaches the distant observer at any finite time.

This is what the astronaut and the distant observer find, by direct construction, with no formal apparatus more sophisticated than ``count the flashes.'' What they find is the constructive content of the metric-singularity theorem. The astronaut's horizon-crossing event is not an event that occurs at any finite distant-observer time. It is, from the distant observer's vantage, an asymptotic limit --- approached but never reached.

But the distant observer is not a special case. The distant observer is every external observer. By the same construction, every external observer finds the same thing: the astronaut's horizon-crossing event is forever in their causal future, asymptotically approached. No external observer at any finite time has the horizon-crossing event in their causal past. The horizon-formation event, the moment at which the collapsing star's surface crosses what the standard reading calls the horizon, is the same kind of event from every external vantage: asymptotic, never finite-time.

This is the third route. It uses no global theorems, no asymptotic structure of null infinity, no formal apparatus beyond Premise 1 (causal locality, which the construction tracks directly) and the standard relativistic redshift (which Premise 4 contains). And it reaches, by an entirely different style of reasoning, the same conclusion the first two routes reached.

\section*{8. The fourth route: four mutually-entailing arguments}

The third article in the programme's popular exposition~\cite{JanzenArticle3} develops the no-horizon-formation conclusion through a fourth route that is, structurally, a tetrahedron of four mutually-entailing arguments. Each face of the tetrahedron is a distinct argument; the six edges express the entailments among them.

\textit{The observability constraint.} The horizon-formation event lies on the boundary of $J^-(\mathscr{I}^+)$, and no signal from it can reach $\mathscr{I}^+$. This is a strict causal-structure consequence, and it forbids the horizon-formation event from being in the causal past of any external observer.

\textit{The worldtube replacement argument.} The horizon, by its global causal definition, depends on the entire future development of the spacetime. As matter continues to collapse, accrete, or otherwise alter the spacetime's future, the location of $H^+$ alters. The horizon is not a stable feature present at any finite exterior time; it is a sequence of would-be locations, each redefined by the next finite-time alteration of the future.

\textit{The asymptotic-nature argument.} The horizon is the limiting null surface at $\mathscr{I}^+$. Treating it as a finite-time structure is a category error about the asymptotic nature of the definition.

\textit{The generator argument.} The local-geometric content of the asymptotic-nature argument: the asymptotic alignment of horizon generators (route 2 above).

Each of the four arguments is an independent route to the same conclusion. Each operates on different terrain. Each requires denying a different one of the minimal premises to defeat. The observability constraint requires denying causal locality (Premise 1); the worldtube replacement requires denying the global form of the horizon definition (Premise 2); the asymptotic-nature argument requires denying the asymptotic structure of $\mathscr{I}^+$ (Premise 3); the generator argument requires denying the standard alignment of horizon generators (Premise 4).

The tetrahedron is the fourth route in the sense that it is its own thing: a methodological-architectural contribution showing that the conclusion is reachable from four distinct angles that mutually entail one another. But it is also, in a sense, the move I am writing this essay to articulate. The four arguments together exhibit the structural-redundancy pattern at the small scale: a single conclusion underwritten by four independent routes.

\section*{9. Four more routes: the algebraic, geometric, structural, and cosmological}

The programme's other four papers extend the constellation through four more routes that operate on terrain quite different from the routes already given.

The \textit{circle paper}~\cite{JanzenCircle} exhibits the Schwarzschild geometry as a circle. This is not metaphor. The Schwarzschild solution, written in a particular radial parametrisation, takes the form of a cycloid: the radial coordinate $r$ varies periodically along an angular parameter $z$, with $r(z) = M(1 + \cos z)$. The standard ``singularity'' at $r = 0$ and the standard ``horizon'' at $r = 2M$ are two non-degenerate critical points of this cycloid. They are of identical analytic type, and not by coincidence: the two critical points are the two poles of a single circle, identical because a circle has no distinguished point --- the conclusion the geometry forces, not one the inspection merely proposes. The horizon-singularity asymmetry of the standard reading --- in which one critical point is classified as a removable coordinate singularity and the other as a true curvature singularity --- is an artefact of the asymmetric classification, not a feature of the geometry. The route is algebraic: the methodological move applied not to the global causal structure of the spacetime but to the algebraic content of the standard radial coordinate.

The \textit{slicing paper}~\cite{JanzenSlicing} unifies Schwarzschild and de Sitter as the two limits of one intrinsic radial curve. The Schwarzschild--de Sitter family is, this paper shows, rigid in the sense that it has no continuous moduli: the family is a set of slicings of one de Sitter manifold, the Schwarzschild and de Sitter forms two limits of one underlying construction and---at the level at which CR individuates---representations of one ontological layer. The standard reading, in which Schwarzschild and de Sitter are two distinct geometries on two distinct manifolds, is incorrect. The route is geometric: the methodological move applied to the geometric structure of the SdS family, with the conclusion that what looks like two autonomous geometries is one structure read two ways.

The \textit{groupoid paper}~\cite{JanzenGroupoid} develops the algebraic content of the slicing paper's groupoid. The discrete morphisms among admissible charting vantages organise into the dihedral group $D_3 \cong S_3$. The Schwarzschild description is the canonical asymmetric realisation, in which a partial arc of the slicing curve is swept about the off-axis point $r=0$ --- a genuine point of the manifold, but not a centre of symmetry. The horizon-versus-singularity asymmetry is, on this reading, a signature of the asymmetric sweep, not a feature of the geometry. The route is algebraic-geometric: the methodological move applied to the structure of admissible descriptions of one geometry. (Together with the slicing paper, the groupoid paper completes a single route: the geometric-algebraic establishment of rigidity plus the sweep-pivot diagnostic.)

The \textit{cosmology paper}~\cite{JanzenCRcosmology} establishes the cosmological completion. A causal reassignment of de Sitter space --- the promotion of a null direction at the horizon to the fundamental timelike congruence --- produces flat $\Lambda$CDM as an algebraic identity. The reassignment is forced rather than postulated: only the Nariai configuration of the SdS family admits a comoving curve along a null generator without pivoting into a horizon. The horizon's null direction (selected asymptotically in collapse) and the cosmic temporal direction (in the Nariai cosmology) are two descriptions of the same intrinsic structure. The route is cosmological: the methodological move applied to the structural identification that ties the no-horizon-formation result to the cosmological completion of the framework.

And one further route, of a different character: the \textit{parity-and-compatibility} result, established in the programme's first two popular-exposition articles. The result is that the framework's reading reproduces every observational fact about exterior regions that the standard reading reproduces. The deviations between the two readings concern only what is taken to exist in finite time \emph{inside} a horizon (which the framework says is nothing). The route is structural: it establishes that the framework is a structural augmentation of general relativity, not a competing theory of gravity --- it makes no observational predictions that conflict with standard GR for any exterior region of any solution. To defeat the framework on this route, an objector would have to identify an exterior-region observational prediction on which the framework and standard GR differ. None has been found.

I have now described all eight routes: the global causal theorem (§5), the asymptotic alignment of null directions (§6), the constructive limit walked through what an observer finds (§7), the Article-3 tetrahedron of four mutually-entailing arguments (§8), and the algebraic, geometric-algebraic, cosmological, and parity-compatibility routes of this section. Each route operates on different terrain, with different load-bearing premises. Each reaches the same conclusion.

A note that may be worth pausing on: I am not claiming that the eight routes are independent in the strong logical sense. They are not. The same minimal premises (causal locality, the global horizon definition, the asymptotic structure of null infinity, the standard structure of null geodesic congruences) are load-bearing in many of them. What is distinct about the routes is not that they assume different things but that they exhibit the conclusion from different angles, in different formalisms, against different counter-arguments. To defeat the conclusion, an objector must defeat each of the eight separately, on its own terrain. A single deflection that works against route 1 does not work against route 2, because routes 1 and 2 operate on different formalisms and against different counter-arguments. The robustness is structural.

\section*{10. Why this matters: the architecture of robustness}

A single careful argument, however well-constructed, can be deflected by a single careful deflection. This is a fact about how scientific arguments are evaluated by the field. A referee who reads the argument can, with effort, find a way to disagree with one of its premises or one of its inferential steps. Whether the disagreement is correct or not, the argument is, in the field's record, deflected: the referee's report exists, the argument's conclusion is contested, and the field's response is to wait for further work that addresses the deflection.

This is not a complaint. It is how scientific verification works in practice. The single-deflection vulnerability is a feature of the system, not a bug. Most arguments that face it are arguments that should be deflected --- arguments that genuinely have a problem the deflection identifies. The field's verification practices are designed to deflect bad arguments. Good arguments are expected to survive deflection by being correct.

But there is a specific class of arguments for which the single-deflection vulnerability becomes a structural problem. These are arguments that bear on foundational questions: questions about what a formalism is taken to describe, about which standardly-accepted features of a theory are actually derived from premises and which are interpretive habits, about whether the field's everyday reading of its central theories is the careful reading. Foundational arguments of this kind are not deflected because they are wrong (in many cases they are correct); they are deflected because the referees applying their training to them are reading the arguments through the very interpretive habits the arguments are correcting. The field's verification practices, when applied to a foundational correction, do not verify; they deflect.

The deflection, in such cases, looks identical in the field's record to the deflection of a bad argument. The referee writes a report; the report identifies what the referee takes to be an error; the journal rejects on the basis of the report. The field's record contains no indication that the deflection was the application of an interpretive habit rather than the identification of a substantive error. The two are indistinguishable from outside.

The structural-redundancy architecture --- the methodological move of establishing the same conclusion along multiple structurally distinct routes --- is the response to this dynamic. A single argument can be deflected by a single referee. Eight arguments operating on eight different terrains, with different load-bearing premises, against different counter-arguments, cannot. A referee who deflects route 1 by importing an alternative definitional framework for $H^+$ has not deflected route 3 (which uses no formal definition of $H^+$ at all, only the constructive content of what the astronaut and the distant observer find). A referee who deflects route 3 by objecting to the magic flashlight has not deflected route 6 (the algebraic identity of the two critical points of the cycloid, which makes no reference to flashlights or observers). A referee who deflects route 6 by disputing the parametrisation has not deflected route 7 (the rigidity of the SdS family, which is established geometrically and does not depend on the cycloid parametrisation).

To defeat the conclusion, the objector must defeat each of the eight separately. This is not impossible in principle. If the conclusion is wrong, eight independent defeats should be findable. But the architecture forces the objector to engage on eight terrains, not one --- and in doing so, it forces the engagement to be substantive on at least seven of them, because the deflection that works on one terrain will not work on the others.

The architectural feature is, in this sense, methodological rather than evidential. It does not make any of the eight arguments stronger as an argument. What it does is make the constellation as a whole more robust against the specific failure mode by which careful arguments on foundational questions are deflected. The robustness is not in the strength of the routes; it is in the multiplicity of them.

\section*{11. The lineage, honestly}

I am tempted, here, to claim historical precedent for the methodological move. The temptation has a recognisable shape: pick foundational corrections the field eventually accepted, cite the canonical names attached to them, and read the structural-redundancy architecture back into their work. The temptation is to say the move is not new, only newly explicit; that the great minds were quietly doing this all along.

I want to refuse this temptation in print, because giving in to it would commit the same failure mode the methodology's diagnostic apparatus identifies. The field's standard lineage stories about its own corrections are unreliable when checked against primary sources. Reaching for canonical names to authorise a present-day move, without doing the primary-source historiographic work to verify what those figures actually did, is the fabrication-via-pattern-matching failure I would otherwise be writing this essay to diagnose.

What I can offer instead is a different kind of historical evidence: a single specific case, primary-source-verified, that bears on whether the field's accepted lineage stories can be trusted in general.

I spent two months in 2014 reading primary-source documents from the Albert Einstein Archives at the Hebrew University of Jerusalem. I never went to Jerusalem. I corresponded by email with the archivist Barbara Wolff, ordered items from her list, paid the access fees by credit card, and read what arrived. The work covered Einstein's cosmological considerations from 1917 through the early 1930s. The paper that emerged from it documents nine findings, each with primary-source receipts, that are at sharp variance with the standard textbook story. To take three of them, in compact form:

The ``greatest blunder'' remark, taught in introductory courses and recited in countless review articles, has no primary-source attestation in Einstein's writings. Its only documented origin is Gamow's posthumously published 1970 autobiography, which alleges a private conversation in which Einstein used a milder phrase (``biggest blunder he ever made in his life''). The ``greatest mistake'' wording can be traced through a propagation chain in the secondary literature: Gribbin 1986, Dainton 2001, Penrose 2005, Kolb 2007, with no link to a primary source where Einstein actually says these words. The standard story has been mythologised in transit.

Einstein actually decided to reject the cosmological constant in May 1923, eight years before the public 1931 announcement. The decisive document is a postcard to Hermann Weyl (Einstein Archives call number 24-081). He held the decision privately until Hubble's 1929 result made it announceable. The standard story that Einstein turned against the constant in reaction to Hubble's discovery reverses the actual chronology: Einstein had already turned, and was waiting for empirical cover.

Hubble's 1929 result did not surprise the theoretical community. Eddington (1923, in the second edition of his \textit{Mathematical Theory of Relativity}) and Weyl (1923, fifth edition of \textit{Raum, Zeit, Materie}) had both publicised the redshift-distance relation as the theoretical expectation from de Sitter's solution. Robertson (1928) had derived the linear relationship to first order. Friedmann's 1922 and 1924 papers, and Lema\^{i}tre's 1927 paper, had developed the expanding-universe solutions before Hubble. The standard narrative --- that cosmologists resisted expansion for a decade before Hubble's empirical confirmation forced their hand --- is not what the documentary record shows. The acceleration in the field's acceptance of expansion around 1930 was driven by the rediscovery of Friedmann's and Lema\^{i}tre's papers by Eddington and de Sitter, not by Hubble's data.

What the case verifies, in compressed form, is that the field's standard lineage stories about a foundational period compress, smooth, and reassign canonical names in ways the documentary record does not support. The actual record shows multiple researchers arriving at compatible results, partially independently, across roughly fifteen years; the institutional response deflecting each result until the deflections became untenable; and canonical naming and standard pedagogical narration arriving substantially after the work, by which point the texture of who-did-what-when had been smoothed away.

This is one corner of one foundational episode in the history of physics. I cannot, on this evidence alone, claim that the structural-redundancy methodology articulated in this essay has been used implicitly throughout the history of foundational corrections. To make that claim with comparable rigor would require primary-source historiographic effort, of the kind I have done in this one period and not in others.

What I can say is that the discipline of the structural-redundancy methodology --- refusing to substitute received memory for primary-source verification, and refusing the canonical-name story when the documents do not support it --- is the same discipline I exhibited reading Einstein's actual postcards. It operates the same way whether you are reading 1923 archival correspondence or evaluating a contemporary claim about no-horizon-formation. The methodology is not retrospectively justified by lineage; it is the discipline that, applied to a question, produces the kind of result a single argument could not. The case for it must stand on the worked example developed in the preceding sections.

The institutional reception problem this essay does not address directly --- why corrections that pass the discipline tend to be received slowly when received at all --- has its own articulation in public-register prose~\cite{JanzenFortress}. There the same underlying problem is named through a different structural framing: a triangulated defence (bad metaphysics, misapplied logic, shallow operationalism / appeal to institutional consensus) maintained by an institutional fortress trained to repel structural challenges through mutually reinforcing patterns. The two framings complement rather than compete. This essay carries the methodological-discipline articulation: what the discipline requires of the practitioner who works on foundational questions. The Fortress essay carries the institutional-diagnostic articulation: what the reception structure does to the practitioner's work when it arrives. Together they describe the same structural situation from the inside and from the outside.

I find I prefer this version of the section to the one I almost wrote.

\section*{12. What the move makes visible}

I have argued that the structural-redundancy architecture is a methodological response to the institutional dynamic by which foundational corrections are deflected. I want to close by noting that the architecture also makes something visible that is otherwise hard to see.

A single argument, even a correct one, can be read as an idiosyncratic position --- one researcher's preferred way of looking at a problem, perhaps interesting, perhaps eccentric, in any case not yet established. The field's training does not by default treat a single careful argument on a foundational question as a settled result. It treats it as an offer, to be evaluated, accepted or deflected.

Eight structurally distinct routes to the same conclusion are not, in the same way, an idiosyncratic position. They are the systematic exhibition of a feature of the formalism. The conclusion is not one researcher's preferred reading; it is what the formalism says under minimal premises, across multiple distinct angles of approach. The architecture makes visible the difference between a thesis being offered and a structural feature being exhibited.

This visibility matters because it changes what is required to engage the work. A thesis can be engaged by disagreement: by saying ``I think you are wrong about $X$,'' and providing reasons. A structural feature of a formalism cannot be engaged by disagreement, only by identification of where the formalism's standard reading and the careful reading diverge. The structural-redundancy architecture forces the engagement to be substantive: not ``I disagree with this conclusion,'' but ``I disagree with this conclusion because at premise $P$ I would substitute $Q$,'' and ``the substitution of $Q$ for $P$ defeats route 1 in this way, route 2 in this way, route 3 in this way, $\ldots$ route 8 in this way.''

That is the level of engagement at which foundational corrections actually settle. The architecture is what forces the engagement to that level.

\section*{13. Leaving the cave}

The prisoners in Plato's allegory are not made to leave the cave by being shown better shadow-science. The shadow-science is, in its own register, sound: the prisoners' taxonomies are accurate, their predictions are reliable, their tests distinguish good shadow-theories from bad ones. What makes the prisoners leave the cave is the recognition that the shadow-science, for all its sophistication, is the study of shadows of something else --- and that the something else is accessible if one is willing to turn around.

The methodological move I have described is, in foundational physics, the turning around. It does not require new formalism. It does not require new instruments. It does not require new observations. It requires the careful reading of the formalism the field already accepts, under premises the field already accepts, with the willingness to follow the conceptual implications wherever they lead.

What makes the move possible is the willingness to step outside the field's interpretive training for long enough to read what the formalism actually says. What makes the move robust is the structural-redundancy architecture: eight structurally distinct routes to the same conclusion, on eight different terrains, against eight different deflections.

The result, in the case of the no-horizon-formation conclusion, is a clean technical correction to a foundational misreading that has propagated through the field for fifty years. The correction does not modify general relativity. It does not propose new physics. It reads the existing physics carefully and reports what the careful reading reveals. The shadows on the wall are real shadows. They are exactly as the shadow-science describes. They are also shadows of something else, and the something else is accessible if one is willing to turn around.

\bigskip

\begin{thebibliography}{9}

\bibitem{JanzenBHcausality}
D.~Janzen,
\emph{The event horizon is a metric singularity: a missing definition in general relativity},
in preparation.

\bibitem{JanzenCircle}
D.~Janzen,
\emph{Schwarzschild as a circle: the maximal extension as analytic continuation through two structurally identical critical points},
in preparation.

\bibitem{JanzenSlicing}
D.~Janzen,
\emph{The Schwarzschild--de Sitter slicing curve: horizons as turning points, and de Sitter and Schwarzschild as the two limits of one construction},
in preparation.

\bibitem{JanzenGroupoid}
D.~Janzen,
\emph{The Schwarzschild--de Sitter description groupoid: generators, relations, and Schwarzschild as the asymmetric realisation},
in preparation.

\bibitem{JanzenCRcosmology}
D.~Janzen,
\emph{Null-Boundary Correspondence and Non-Synchronous Cosmology in a Layered Geometric Framework},
in preparation.

\bibitem{JanzenFramework}
D.~Janzen,
\emph{Cosmological Relativity: a structural augmentation of general relativity},
in preparation.

\bibitem{JanzenArticle3}
D.~Janzen,
\emph{The tetrahedral proof: four mutually-entailing arguments that black holes do not form in finite time},
cosmiCave essay (2026).

\bibitem{JanzenFortress}
D.~Janzen,
\emph{The Fortress of Misunderstanding Physics},
cosmiCave essay (August 2025).

\end{thebibliography}

\end{document}
